\edef\refsec{\refnumsec}
\edef\refeq{\refnumeq}
\edef\refalg{\refnumalg}


\section{Numerical Solution}\label{\refsec}
This section documents how the model is solved numerically. Three ideas organise the approach, and they recur throughout:
\begin{enumerate}
    \item \emph{Policies are identified from first order conditions, solved as root problems.} Combined with the result of \cref{\refmodelsec:PEE} that the predetermined savings shares $s_{t-1,i}/s_{t-1}$ are closed-form functions of the current tax, this keeps the entire distribution of past savings out of the political state space. In this model variant the reduction goes furthest: under log preferences the political problem carries no endogenous state at all and \emph{decouples across periods} into $T$ independent scalar problems (\cref{\refsec:PEELOG}); under CRRA preferences the aggregate savings level returns as the single state (\cref{\refsec:PEE}).
    \item \emph{Root problems are solved on grids that handle corner solutions and multiple roots by construction.} A bounded reparameterization lets gradient-based solvers respect the admissible policy interval, corners are detected rather than assumed away, and the same grid evaluations are reused both to reconstruct the political objective --- which is what selects among several candidate optima --- and to approximate the derivatives that have no closed form. \Cref{\refsec:RobustRoot} collects this machinery once; the subsections after it apply it.
    \item \emph{Choices whose interior status is itself the question are solved on the objective, not the first order condition.} The endogenous choice of the system design $\theta$ (\cref{\refsec:ESC}) is layered on top of the tax machinery by evaluating the political objective on a grid of designs and maximising directly, so that a corner design is \emph{observed} rather than inferred from a solver's failure to converge --- the central results of that exercise concern precisely whether the choice is interior.
\end{enumerate}
A complete implementation accompanies the paper in the online repository at \url{https://github.com/ChampionApe/SSDclaude}: the model-specific code is in \texttt{python/US} and the generic grid-search routines in \texttt{python/gridsearch}. The presentation here is self-contained and does not presume familiarity with the code; conversely, the docstrings in the code cite the equation labels used below.

Before we outline the specific solutions, we collect all relevant baseline equations in one place with explicit reference to the arguments that may change during a solution (all of these are simply copies from \cref{\refmodelsec}, but collected here for simplicity). We reference them with explicit arguments for the parts that may change during solution/simulation/calibration stages:
\begin{subequations}\label{\refeq:auxiliary}
    \begin{align}
        \Gamma_{s,t}\left(\bm{B}_{t+1}, \tau_{t+1}, \theta_{t+1} \right) &= \frac{1}{1+\xi}\frac{\sum_i \frac{\gamma_{t,i}\eta_{t,i}^{1+\xi}}{X_{t,i}^{\xi}}\frac{B_{t+1}^i}{1+B_{t+1}^i}}{1+\frac{1-\alpha}{\alpha}\tau_{t+1}\left(\theta_{t+1}+(1-\theta_{t+1})\sum_i\frac{\gamma_{t,i}}{1+B_{t+1}^i}\right)},
        \label{\refeq:auxiliary:Gammas} \\
        \Theta_{h,t}\left(\tau_t, \tau_{t+1}, \theta_{t+1}, \Gamma_{s,t}\right) &= \left(\Gamma_{h,t}\right)^{\frac{1+\xi}{1+\alpha\xi}}\left(\frac{(1-\alpha)(1-\tau_t)}{\Gamma_{h,t}-\frac{1-\alpha}{\alpha}\theta_{t+1}\tau_{t+1}\Gamma_{s,t}}\right)^{\frac{\xi}{1+\alpha\xi}}
        \label{\refeq:auxiliary:Thetah} \\
        \Theta_{h,T}\left(\tau_T\right) &= \left(\Gamma_{h,T}\right)^{\frac{1}{1+\alpha\xi}}\left((1-\alpha)(1-\tau_T)\right)^{\frac{\xi}{1+\alpha\xi}}
        \label{\refeq:auxiliary:ThetahT} \\
        \Theta_{s,t}\left(\Theta_{h,t}, \Gamma_{s,t}\right) &= \left(\frac{\Theta_{h,t}}{\Gamma_{h,t}}\right)^{\frac{1+\xi}{\xi}} \Gamma_{s,t}
        \label{\refeq:auxiliary:Thetas} \\
        s_t(\Theta_{s,t}, s_{t-1}) &= \Theta_{s,t} \left(\frac{s_{t-1}}{\nu_t}\right)^{\frac{\alpha(1+\xi)}{1+\alpha\xi}}
        \label{\refeq:auxiliary:s} \\
        h_t\left(\Theta_{h,t}, s_{t-1}\right) &= \Theta_{h,t}\left(\frac{s_{t-1}}{\nu_t}\right)^{\frac{\alpha\xi}{1+\alpha\xi}} \label{\refeq:auxiliary:h} \\
        B_{t+1}^i\left(s_t, h_{t+1}\right) &=\left(\beta_{t,i}\right)^{\rho}\left[\frac{\alpha}{p_t}\left(\frac{s_t}{\nu_{t+1}h_{t+1}}\right)^{\alpha-1}\right]^{\rho-1}
        \label{\refeq:auxiliary:B} \\
        R_{t+1}\left(s_t, h_{t+1}\right) &= \alpha \left(\frac{s_t}{\nu_{t+1}h_{t+1}}\right)^{\alpha-1} \label{\refeq:auxiliary:R} \\
        s_t\left(h_t, \Gamma_{s,t}\right) &= \left(\frac{h_t}{\Gamma_{h,t}}\right)^{\frac{1+\xi}{\xi}}\Gamma_{s,t}
        \label{\refeq:auxiliary:sFromH} \\
        \frac{s_{t,i}}{s_t}\Bigg(\bm{B}_{t+1}, \tau_{t+1}, \theta_{t+1}\Bigg) &=\frac{\frac{B_{t+1}^i}{1+B_{t+1}^i}\frac{\eta_{t,i}^{1+\xi}}{X_{t,i}^{\xi}}}{\sum_{j} \frac{\gamma_{t,j}\eta_{t,j}^{1+\xi}}{X_{t,j}^{\xi}}\frac{B_{t+1}^j}{1+B_{t+1}^j}}\left[1+\frac{1-\alpha}{\alpha}\tau_{t+1}\left(\theta_{t+1}+(1-\theta_{t+1})\sum_{j}\frac{\gamma_{t,j}}{1+B_{t+1}^j}\right)\right] \notag \\
        -&\frac{1}{1+B_{t+1}^i}\frac{1-\alpha}{\alpha}(1-\theta_{t+1})\tau_{t+1} - \frac{\eta_{t,i}^{1+\xi}/X_{t,i}^{\xi}}{\Gamma_{h,t}}\frac{1-\alpha}{\alpha}\theta_{t+1}\tau_{t+1} \label{\refeq:auxiliary:si_s}
    \end{align}
\end{subequations}
Two conventions are worth flagging before we use these. First, \eqref{\refeq:auxiliary:si_s} is written with the vintage $t$; the political problem at $t$ needs it one period back, i.e.\ evaluated at $\bm{B}_t$, $\tau_t$, $\theta_t$ and the $t-1$ parameter vintage, and it is easy to introduce an off-by-one error here. Second, the model carries no informal type, so the coefficient $\kappa_t$ appearing in the informal-sector documentation has collapsed to $p_t$ and cancels against the $p_t$ in the discounted pension term; the group $\frac{1-\alpha}{\alpha}\tau_{t+1}$ above is what remains of it.

\subsection{The economic equilibrium}\label{\refsec:EE}

\smalltitle{With log preferences.}
Assume that we know the path of taxes $\tau_t$, pension characteristics $\theta_t$, and an initial exogenous level of savings $s_0$. In this case, recall that we simply have $B_{t+1}^i = \beta_{t,i}$. Given this, we can solve for the economic equilibrium without any numerical root finding:
\begin{itemize}
    \item First compute $\Gamma_{s,t}$ from \eqref{\refeq:auxiliary:Gammas} (only defined for $t<T$).
    \item Then compute finite horizon $\Theta_{h,t}$ from stacking \eqref{\refeq:auxiliary:Thetah} for $t<T$ and \eqref{\refeq:auxiliary:ThetahT} for $t=T$.
    \item Then compute $\Theta_{s,t}$ from \eqref{\refeq:auxiliary:Thetas} (only defined for $t<T$ as well).
    \item Given $s_0$ and $\Theta_{s,t}$ we can iterate forward on \eqref{\refeq:auxiliary:s} to solve for $s_t$ for all $t<T$. For $t=T$ we have $s_T = 0$ per construction.
\end{itemize}
As explained in \cref{\refmodelsec}, this set of core variables allows us to solve for individual labour supply and savings, factor prices, pension benefits, and finally consumption levels (from budgets).


\smalltitle{With CRRA preferences.}
Assume again that we know the path of taxes $\tau_t$, pension characteristics $\theta_t$, and an initial exogenous level of savings $s_0$. The core of the economic equilibrium can then be reduced to identifying vectors of $\Gamma_{s,t}, h_t, s_t$. The three vectors are identified by three sets of conditions (we have a square system of nonlinear equations):
\begin{itemize}
    \item The finite horizon labour supply condition requires \eqref{\refeq:auxiliary:h} with $\Theta_{h,t}$ being computed using \eqref{\refeq:auxiliary:Thetah} for $t<T$ and \eqref{\refeq:auxiliary:ThetahT} for $t=T$.
    \item The finite horizon savings decision requires solving for $t<T$ (with $s_T = 0$); we use \eqref{\refeq:auxiliary:sFromH} for this.
    \item Finally, $\Gamma_{s,t}$ is also only defined for $t<T$; we use \eqref{\refeq:auxiliary:Gammas} for this, evaluated with $B_{t+1}^i(s_t, h_{t+1})$ from \eqref{\refeq:auxiliary:B}.
\end{itemize}
Given a solution to this core of the economic equilibrium ($s_t, h_t, \Gamma_{s,t}$), the rest of the outcomes can be solved as in the log case.


\smalltitle{Initializing with steady state savings.}
When we solve for the economic equilibrium given vectors of $\tau_t$ and $\theta_t$, our default option for the initial exogenous level of savings $s_0$ is to solve for steady state levels using parameter values from the first period ($t=1$). This default is also what makes the scale invariance of \eqref{\refmodeleq:scaleInvariance} operative: $s_0$ is then a function of the same parameters as everything else and rescales with them, which it would not do if imposed exogenously.

With log preferences, we have that $B^i = \beta_i$ and then we can solve for $\Gamma_{s}$ directly in closed form using the general equation \eqref{\refeq:auxiliary:Gammas}. Given this, we can then compute steady state savings from \eqref{\refeq:steadystate_LOG:s}:
\begin{subequations}\label{\refeq:steadystate_LOG}
    \begin{align}
        \Gamma_s =& \frac{1}{1+\xi}\frac{\sum_i \frac{\gamma_{i}\eta_{i}^{1+\xi}}{X_{i}^{\xi}}\frac{B^i}{1+B^i}}{1+\frac{1-\alpha}{\alpha}\tau\left(\theta+(1-\theta)\sum_i\frac{\gamma_{i}}{1+B^i}\right)} \label{\refeq:steadystate_LOG:Gammas}\\
         s^* =& \Gamma_h^{1+\xi}\left[\left(\frac{(1-\alpha)(1-\tau)}{\Gamma_h-\frac{1-\alpha}{\alpha}\theta\tau\Gamma_s}\right)^{1+\xi}\frac{\Gamma_s^{1+\alpha\xi}}{\nu^{\alpha(1+\xi)}}\right]^{\frac{1}{1-\alpha}} \label{\refeq:steadystate_LOG:s}
    \end{align}
\end{subequations}

With CRRA preferences, we have to identify $s^*$ numerically instead. This can technically be done in several ways, but we prefer the following as it allows us to solve the problem as a bounded scalar search where there are robust and fast algorithms available (e.g.\ golden-section search). Specifically, we solve for the $\Gamma_s$ that satisfies condition \eqref{\refeq:steadystate_CRRA:Gammas}:
\begin{subequations}\label{\refeq:steadystate_CRRA}
\begin{align}
    &\Gamma_s(1+\xi)\left[1+\frac{1-\alpha}{\alpha}\tau\left(\theta+(1-\theta)\sum_i\frac{\gamma_i}{1+B^i(\Gamma_s)}\right)\right] = \sum_i\frac{\gamma_i\eta_i^{1+\xi}}{X_i^{\xi}} \frac{B^i(\Gamma_s)}{1+B^i(\Gamma_s)}, \label{\refeq:steadystate_CRRA:Gammas} \\
    &\text{where }B^i(\Gamma_s) = \left(\beta_i\right)^{\rho} \left(\frac{\alpha}{p}\frac{\Gamma_h-\frac{1-\alpha}{\alpha}\theta\tau\Gamma_s}{(1-\alpha)(1-\tau)}\frac{\nu}{\Gamma_s}   \right)^{\rho-1} \label{\refeq:steadystate_CRRA:Bi}
\end{align}
\end{subequations}
Two bounds on $\Gamma_s$ matter, and only the first is the obvious one. The right hand side of \eqref{\refeq:steadystate_CRRA:Gammas} is at most $\Gamma_h$ and the bracket on the left is at least one, so with constant $B^i$ across types
\begin{align}\label{\refeq:steadystate_CRRA:bounds}
    \Gamma_s \in \left(0, \frac{\Gamma_h}{1+\xi}\frac{B}{1+B}\right).
\end{align}
This is a bound on where the root \emph{is}. The second is a bound on where the residual can be \emph{evaluated at all}. Both $\Theta_{h,t}$ in \eqref{\refeq:auxiliary:Thetah} and $B^i$ in \eqref{\refeq:steadystate_CRRA:Bi} carry the factor $\Gamma_h-\frac{1-\alpha}{\alpha}\theta\tau\Gamma_s$, which vanishes at
\begin{align}\label{\refeq:steadystate_CRRA:cap}
    \Gamma_s^{\max}\left(\tau,\theta\right) = \frac{\Gamma_h\,\alpha}{(1-\alpha)\,\theta\tau},
\end{align}
and is negative above it. A steady state always satisfies $\Gamma_s<\Gamma_s^{\max}$ --- $\Theta_{h}$ must be real and positive --- so \eqref{\refeq:steadystate_CRRA:cap} never excludes the root. What it does exclude is a chunk of the search interval: raising $\Gamma_s$ past $\Gamma_s^{\max}$ raises a negative number to the fractional power $\rho-1$, and the residual returns NaN rather than a finite value of the wrong sign. A bracketing method then fails on the NaN, reporting a problem with the solver rather than with the interval it was handed.

The practical bracket is therefore $\left(0,\min\left\lbrace u_{\Gamma}, \varsigma\,\Gamma_s^{\max}(\tau,\theta)\right\rbrace\right)$ for a fixed $u_{\Gamma}$ and a safety factor $\varsigma$ slightly below one. It is worth being explicit about why a constant $u_{\Gamma}$ alone will not do, because it \emph{appears} to work: with $\Gamma_h = 1$, \eqref{\refeq:steadystate_CRRA:cap} is $\alpha/((1-\alpha)\theta\tau)$, which at $\alpha = 0.43$ and the Argentina calibration stays above a constant such as $0.75$ over the whole unit interval of $\tau$, while at $\alpha = 0.30$ it falls to $\approx 0.58$ as $\tau\rightarrow 1$. A constant that is safe in one calibration is thus not safe in another, and since the steady state is evaluated at trial taxes spanning $[l,u]$ while solving for the politico-economic equilibrium, the infeasible region is genuinely visited rather than merely nearby. Retuning the constant would only move the $\tau$ at which the problem reappears; tying the bracket to \eqref{\refeq:steadystate_CRRA:cap} removes it. Note finally that $\Gamma_h$ is kept explicit in both \eqref{\refeq:steadystate_CRRA:bounds} and \eqref{\refeq:steadystate_CRRA:cap} even though both model variants normalize it to one, so that the bounds remain consequences of the model rather than of the normalization.

\subsection{Robust root finding with bounds}\label{\refsec:RobustRoot}
In the following subsections, we identify policies that maximize the political objective function by solving the first order condition as a root-finding problem instead of using the political objective function directly. This is paramount to avoid having the entire distribution of past savings as state variables for the political problem. Also, we will generally only work with $\tau \in[0,1]$ for numerical stability. To deal with these two points (that we work with first order conditions and only want to evaluate most functions on $\tau\in[0,1]$), we set up an auxiliary root-finding problem.

Notationally, let $z(\tau) = f\left([\tau]_{[0,1]}\right)$ denote the first order condition evaluated in the interior $[\tau]_{[0,1]} \equiv \min(1, \max(0,\tau))$. Instead of trying to solve for $f(\tau) = 0$, we suggest solving for
\begin{align}\label{\refeq:root}
    h(z,\tau) = z - |z| \underbrace{\left[a_0\min\left(\tau,0\right) + a_1\left(\max(\tau,1)-1\right)\right]}_{\equiv g(\tau)} = 0,
\end{align}
where $a_0, a_1>0$ are some technical parameters. In this case we have three regions:
\begin{align*}
    h\left(z,\tau\right) = \left\lbrace \begin{array}{ll} f(0)-|f(0)| a_0 \tau, & \tau<0 \\
    f(\tau), & \tau\in[0,1] \\
    f(1)-|f(1)|a_1(\tau-1), & \tau>1\end{array}\right., && \dfrac{\partial h}{\partial \tau} = \left\lbrace \begin{array}{ll} -|f(0)|a_0, & \tau<0 \\
    \dfrac{\partial f}{\partial \tau}, & \tau\in[0,1] \\
    -|f(1)|a_1, & \tau>1\end{array}\right.
\end{align*}

If the function $f(\tau)$ is well behaved (i.e.\ identifies a maximum), we can solve simply using $h(z(\tilde{\tau}), \tilde{\tau}) = 0$ and then report the bounded solution as $\tau = [\tilde{\tau}]_{[0,1]}$. As long as the function is well behaved, we can use gradient-based solvers to identify roots.

If the function is not well behaved, we may still use $h(z, \tau)$ to detect boundary solutions in some instances. Assume for instance that $f(\tau) >0$ for all $\tau\in[0,1]$, but that $\partial f/\partial \tau <0$. In this case a gradient-based algorithm would have a hard time identifying a solution starting in the interior, as it would suggest trying out lower values of $\tau$. However, note that for $\tau\geq 1$, the method still works and identifies $\tilde{\tau} = 1+1/a_1$ and thus the boundary solution of $\tau = 1$. Similarly, assume that $f(\tau)<0$ for all $\tau\in[0,1]$ and that $\partial f/\partial \tau >0$. In this case, the same problem exists for a gradient-based solver in the interior, but searching outside the bounds still ensures that we identify the correct root $\tilde{\tau} = -1/a_0$ and thus $\tau = 0$.

\subsection{The politico-economic equilibrium with log}\label{\refsec:PEELOG}
The politico-economic equilibrium is straightforward to solve with log preferences. Once again, the problem could be set up and solved in many different ways. In the following, we go through a couple of options that trade off speed with robustness. We refer to these as ``algorithms'' here to have a clear reference, but it may be more suitable to think of them as solution approaches.

Throughout, let $z_t$ denote the marginal effect of policy on the political objective,
\begin{align}\label{\refeq:LOG:foc}
    z_t &= \omega\sum_{i}\left(\gamma_{t-1,i}p_{t-1}\mu_{t-1,i} \frac{d\upsilon_{2,t}^i}{d\tau_t}\right)+\nu_t\sum_i\left(\gamma_{t,i}\mu_{t,i} \frac{d\upsilon_{1,t}^i}{d\tau_t}\right),
\end{align}
where the marginal effects on indirect utilities are defined by \eqref{\refmodeleq:PEELOG} for $t<T$ and \eqref{\refmodeleq:terminalPEELOG} for $t=T$.

\smalltitle{The problem separates into $T$ scalar problems.}
As established in \eqref{\refmodeleq:PEELOG:decoupling}, $z_t$ is a function of $\tau_t$ alone. This is the single most important structural fact for the numerical solution, and it is specific to the model without informal households. Three consequences follow. First, there is no recursion: the periods may be solved in any order, or all at once. Second, the $T$-dimensional root problem $\bm{z}(\bm{\tau}) = \bm{0}$ has a \emph{diagonal} Jacobian, so a vectorized Newton step is nothing but $T$ independent scalar Newton steps sharing one function evaluation --- there is no cross-period information for a solver to exploit, and none to lose by ignoring it. Third, robustness is local: a period in which the first order condition is badly behaved cannot contaminate any other period, which is not true of the backward recursions used for the models with an informal sector.

Everything in this subsection could therefore be stated for a single, arbitrary $t$. We keep the vector notation only because the implementation evaluates all periods in one pass.

\begin{alg}[Gradient-based roots for the politico-economic equilibrium with log preferences]
\label{\refalg:LOG:fast}
This is barely an algorithm, in the sense that we do not include fallback options, but simply presume that the first order condition is well behaved. Let $\tilde{\bm{\tau}} \in \mathbb{R}^{T}$ be a candidate vector of policies and $\bm{\tau}$ the bounded version where each $\tau_t \in [l, u]$ and where $l\geq 0$ and $u\leq 1$ are stable boundaries within the unit interval where it is safe to evaluate all relevant functions. We solve for the vector by identifying the $\tilde{\bm{\tau}}$ that solves \eqref{\refeq:root} period by period, and then report $\bm{\tau}$ as the bounded version.
\end{alg}

\smalltitle{Behaviour at the boundaries.}
Both $d\upsilon_{1,t}^i/d\tau_t$ and $\partial \ln(\Theta_{h,t})/\partial\tau_t$ carry a factor $1/(1-\tau_t)$, so $z_t\rightarrow -\infty$ as $\tau_t\rightarrow 1$. We therefore always evaluate on $[l,u]$ with $u<1$ strictly. The divergence makes it easy to choose $u$ close enough to $1$ that $z_t(u)<0$, and we assume this throughout. Two consequences follow: the upper corner is then never selected, and $z_t(l)>0$ guarantees by continuity that at least one interior downward crossing exists. The substantive difficulty is thus not existence, but (i) detecting a lower-corner solution and (ii) choosing between several interior crossings when they occur. Both are worth verifying numerically rather than assuming, since $z_t(u)<0$ is a property of the calibration, not an identity.

\smalltitle{The extended grid.}
Let $\mathcal{T} = \left\lbrace \tau^{(1)},...,\tau^{(M)}\right\rbrace$ with $\tau^{(1)} = l$ and $\tau^{(M)} = u$ denote a grid on the interior, and let the \emph{extended} grid $\tilde{\mathcal{T}}$ add two outer nodes
\begin{align}\label{\refeq:extendedGrid}
    \tilde{\tau}^{(0)} = l-\frac{1}{a_0}-\delta, \qquad \tilde{\tau}^{(M+1)} = u+\frac{1}{a_1}+\delta,
\end{align}
for a small $\delta>0$. The outer nodes sit at (a hair beyond) the two boundary solutions of \eqref{\refeq:root}. Evaluating $h$ at $\tilde{\tau} = l-1/a_0$ gives $h = f(l)+\left|f(l)\right|$, which is identically zero whenever $f(l)<0$; at $\tilde{\tau} = u+1/a_1$ it gives $h = f(u)-\left|f(u)\right|$, identically zero whenever $f(u)>0$. A corner solution is thus encoded as an \emph{exact} zero at an outer node rather than as a sign change, and a pure sign-change test would not detect it. The offset $\delta$ shifts the node just outside the root so that the outermost cell does exhibit a sign change; since $h$ is affine in $\tilde\tau$ outside $[l,u]$, linear interpolation on that cell returns $\tilde\tau = l-1/a_0$ (resp.\ $u+1/a_1$) exactly, and hence $\tau = l$ (resp.\ $\tau = u$) after bounding. Note that $a_0, a_1$ thereby serve double duty --- as the penalty slopes in \eqref{\refeq:root} and as the placement of the outer nodes --- and the two uses must be kept consistent.

\smalltitle{Selecting among several roots.}
Solving $z_t = 0$ is only a necessary condition for the political optimum: an upward crossing of $z_t$ is a local \emph{minimum} of $\mathcal{W}_t$, and when several downward crossings exist the first one located need not be the global maximum. Since a grid search delivers $z_t$ everywhere on $\mathcal{T}$ in any case, we can restore the maximization without any additional evaluations of $z_t$.

Let $\hat{z}_t$ denote the piecewise-linear interpolant of $z_t$ through the nodes of $\mathcal{T}$ --- the same interpolant that is used to locate crossings --- and define
\begin{align}\label{\refeq:objectiveProfile}
    \hat{\mathcal{V}}_t(\tau) = \int_l^{\tau}\hat{z}_t(x)\mbox{ }dx.
\end{align}
Because $z_t = d\mathcal{W}_t/d\tau_t$, $\hat{\mathcal{V}}_t$ approximates the political objective up to an additive constant. It is piecewise quadratic and continuously differentiable with $\hat{\mathcal{V}}_t^{\prime} = \hat{z}_t$, so its interior local maxima are exactly the downward crossings of $\hat{z}_t$, and
\begin{align}\label{\refeq:candidates}
    \arg\max_{\tau\in[l,u]}\hat{\mathcal{V}}_t(\tau)\in\mathcal{C}_t \equiv \left\lbrace l,u \right\rbrace\cup\left\lbrace \tau\in(l,u)\mbox{ }:\mbox{ }\hat{z}_t(\tau) = 0\text{ and }\hat{z}_t\text{ crosses from above}\right\rbrace.
\end{align}
The candidate set $\mathcal{C}_t$ is finite and cheap to construct, and $\hat{\mathcal{V}}_t$ is evaluated on it \emph{exactly}: at the nodes by the trapezoidal rule,
\begin{align*}
    \hat{\mathcal{V}}_t\left(\tau^{(k+1)}\right) = \hat{\mathcal{V}}_t\left(\tau^{(k)}\right)+\frac{1}{2}\left(z_t\left(\tau^{(k)}\right)+z_t\left(\tau^{(k+1)}\right)\right)\left(\tau^{(k+1)}-\tau^{(k)}\right),
\end{align*}
and, for a crossing $\tau^{\star}$ located inside the cell $\left[\tau^{(k)},\tau^{(k+1)}\right]$, by adding the triangle $\frac{1}{2}z_t\left(\tau^{(k)}\right)\left(\tau^{\star}-\tau^{(k)}\right)$. Setting $\tau_t = \arg\max_{\tau\in\mathcal{C}_t}\hat{\mathcal{V}}_t(\tau)$ resolves multiplicity and corner solutions under a single criterion: $l$ and $u$ compete with the interior crossings on the same footing, so the rule subsumes the boundary check rather than supplementing it. The accuracy demanded of \eqref{\refeq:objectiveProfile} is only that it rank well-separated candidates correctly, which is a far weaker requirement than the accuracy demanded of the located roots themselves.

\begin{alg}[Grid search over the full path]\label{\refalg:LOG:gridsearch}
Fix $\bm{\theta}$ and a global grid $\mathcal{T}$, extended to $\tilde{\mathcal{T}}$ as in \eqref{\refeq:extendedGrid}. Exploiting \eqref{\refmodeleq:PEELOG:decoupling}, all periods are solved simultaneously.

\smallskip\noindent\emph{Evaluate.} Form $z_t$ on $\left\lbrace 1,...,T\right\rbrace\times\tilde{\mathcal{T}}$ in a single vectorized pass, using \eqref{\refmodeleq:terminalPEELOG} in the row $t=T$ and \eqref{\refmodeleq:PEELOG} elsewhere. Every entry is a closed-form expression of $\tau_t$ and the period's parameter vintage; nothing is carried between rows.

\smallskip\noindent\emph{Select.} For each $t$ independently: if $z_t<0$ throughout $\mathcal{T}$, set $\tau_t = l$; if $z_t>0$ throughout, set $\tau_t = u$; otherwise construct $\mathcal{C}_t$ from \eqref{\refeq:candidates} and set $\tau_t = \arg\max_{\mathcal{C}_t}\hat{\mathcal{V}}_t$.

\smallskip\noindent\emph{Polish.} Optionally refine the selected path with Algorithm \ref{\refalg:LOG:fast}, started at the $\tilde{\bm{\tau}}$ implied by $\bm{\tau}$. This is the recommended default, using the grid search to locate the correct branch and the gradient step to attain machine precision.

The cost is $O(T\cdot M)$ evaluations of closed-form expressions, with no derivatives and no initial guess. Since the policy path is smooth in $t$ through the parameter vintages, a windowed variant that evaluates only nodes near the neighbouring periods' solutions is possible and cheaper; unlike in the models with an informal sector, however, the window is a convenience rather than a requirement, and widening it cannot change the answer in any other period.
\end{alg}

\subsection{The politico-economic equilibrium with CRRA}\label{\refsec:PEE}
With CRRA preferences, the past savings level $s_{t-1}$ becomes a relevant state variable for the political problem, and with it the decoupling of \eqref{\refmodeleq:PEELOG:decoupling} is lost: the continuation policy re-enters through $\bm{B}_{t+1}$ and $\Gamma_{s,t}$, and we are back to a backward recursion over policy \emph{functions} $\tau_t(s_{t-1})$. The state-space reduction survives otherwise intact --- the savings \emph{shares} $s_{t-1,i}/s_{t-1}$ remain closed form in $\tau_t$, so the single aggregate level is the only state --- and the machinery of \cref{\refsec:RobustRoot} applies along the policy dimension unchanged. What is new is the grid structure over the state, and a fixed-point problem that pins the endogenous state down.

\smalltitle{Grids.} Let $\mathcal{S} = \left\lbrace s^{(1)}, ..., s^{(M_s)} \right\rbrace$ with $s^{(1)} = l_s>0$ and $s^{(M_s)} = u_s>l_s$ denote a grid of states. Realistic ranges come from steady state dynamics: $s^{\star}(\tau)$ is monotonically decreasing in $\tau$ with $s^{\star} \rightarrow 0$ as $\tau \rightarrow 1$, so $l_s$ is a low value serving numerical stability and the upper bound can be taken as e.g.\ $u_s = 1.25 \times s^{\star}(0)$. Note that $\mathcal{S}$ is a grid in \emph{levels}, so by \eqref{\refmodeleq:scaleInvariance} its placement --- unlike the solution it supports --- depends on the units convention chosen in the calibration and has to be rebuilt, not reused, when $\Gamma_h$ changes. Furthermore, let $\mathcal{S}\tilde{\mathcal{T}}$ denote the Cartesian grid of $\mathcal{S}$ and the extended policy grid $\tilde{\mathcal{T}}$.

For $t<T$ we additionally need a grid of \emph{candidate} values for the endogenous state $s_t$ itself; call it $\mathcal{S}^{\prime}$. It plays a different role from $\mathcal{S}$ --- $\mathcal{S}$ indexes the predetermined states we solve \emph{at}, while $\mathcal{S}^{\prime}$ is searched \emph{over} to locate an equilibrium --- so although the two may hold the same values, it is worth keeping them separate, not least because the accuracy of the search below is governed by $\mathcal{S}^{\prime}$ alone and may warrant a finer spacing.

\smalltitle{The economic equilibrium at $t<T$ is a fixed point in $s_t$.}
Suppose we have solved period $t+1$ and hold continuation policies $\tau^{t+1}(\cdot)$ and $h^{t+1}(\cdot)$, both functions of the state $s_t$ entering $t+1$. For a candidate current tax $\tau_t$ and a candidate state $s_t$, the auxiliary equations \eqref{\refeq:auxiliary} can be evaluated in one forward pass:
\begin{subequations}\label{\refeq:stateApprox}
\begin{align}
    \bm{B}_{t+1} &= \bm{B}_{t+1}\left(s_t,\, h^{t+1}(s_t)\right), \label{\refeq:stateApprox:B}\\
    \Gamma_{s,t} &= \Gamma_{s,t}\left(\bm{B}_{t+1},\, \tau^{t+1}(s_t),\, \theta_{t+1}\right), \label{\refeq:stateApprox:Gammas}\\
    \Theta_{h,t} &= \Theta_{h,t}\left(\tau_t,\, \tau^{t+1}(s_t),\, \theta_{t+1},\, \Gamma_{s,t}\right), \label{\refeq:stateApprox:Thetah}\\
    \Theta_{s,t} &= \Theta_{s,t}\left(\Theta_{h,t},\, \Gamma_{s,t}\right). \label{\refeq:stateApprox:Thetas}
\end{align}
\end{subequations}
The candidate $s_t$ is an equilibrium at $s_{t-1}$ precisely when it reproduces itself through \eqref{\refeq:auxiliary:s}. Defining the residual
\begin{align}\label{\refeq:stateResidual}
    \mathcal{R}_t\left(\tau_t, s_t;\, s_{t-1}\right) \equiv \Theta_{s,t}\left(\tau_t, s_t\right)\left(\frac{s_{t-1}}{\nu_t}\right)^{\frac{\alpha(1+\xi)}{1+\alpha\xi}}-s_t,
\end{align}
the state policy function $s_t(\tau_t, s_{t-1})$ is the root of \eqref{\refeq:stateResidual} in its second argument. This is a genuine root-finding problem in $s_t$ --- not a recursion we can unroll --- because $s_t$ enters the right hand side through $\bm{B}_{t+1}$ and the continuation policies. Note that it is a \emph{root} problem, not a maximisation: any crossing is a solution, so the direction-filtering and objective-integration apparatus of \eqref{\refeq:candidates} is not used here.

\smalltitle{The fixed point is two-dimensional, not three.}
Read \eqref{\refeq:stateApprox} again and note that \emph{none} of $\bm{B}_{t+1}$, $\Gamma_{s,t}$, $\Theta_{h,t}$, $\Theta_{s,t}$ depends on $s_{t-1}$: the predetermined state enters \eqref{\refeq:stateResidual} only through the explicit factor $(s_{t-1}/\nu_t)^{\alpha(1+\xi)/(1+\alpha\xi)}$. The costly part of the state approximation is therefore a function of $(\tau_t, s_t)$ alone, evaluated once on $\tilde{\mathcal{T}}\times\mathcal{S}^{\prime}$, after which the residual for \emph{every} $s_{t-1}\in\mathcal{S}$ follows by multiplying $\Theta_{s,t}$ by a scalar and subtracting $s_t$. What looks like a three-dimensional problem thus costs a two-dimensional one plus a broadcast, which is what keeps the approach affordable as $\mathcal{S}$ is refined.

\smalltitle{Feasibility.}
The root of \eqref{\refeq:stateResidual} need not lie inside $\mathcal{S}^{\prime}$. For a given $s_{t-1}$, large $\tau_t$ depresses savings enough that the implied $s_t$ falls below $l_s$, and conversely at small $\tau_t$; outside that window no equilibrium is located and every object downstream of $s_t$ is undefined. We therefore carry an explicit feasibility mask over $\tilde{\mathcal{T}}\times\mathcal{S}$, marking the pairs $(\tau_t, s_{t-1})$ for which a root was found, and require at least two feasible $\tau_t$ per state before attempting to select a maximum there. Since $s_t(\tau_t)$ is decreasing in $\tau_t$, the feasible set is expected to be a contiguous interval in $\tau_t$, but nothing relies on that: masking infeasible cells individually is no harder and does not silently mis-handle a non-monotone case.

\smalltitle{Which derivatives may be taken on the grid.}
Once $s_t(\tau_t,s_{t-1})$ is known, both marginal indirect utilities take the form (level)$^{1-1/\rho}\times$(log-derivative):
\begin{subequations}\label{\refeq:PEEcrraTerms}
\begin{align}
    \frac{d\upsilon_{1,t}^i}{d\tau_t} &= \left(\hat{c}_{1,t}^i\right)^{1-1/\rho}\frac{d\ln\left(\hat{c}_{1,t}^i\right)}{d\tau_t}, &
    \frac{d\upsilon_{2,t}^i}{d\tau_t} &= \left(c_{2,t}^i\right)^{1-1/\rho}\frac{d\ln\left(c_{2,t}^i\right)}{d\tau_t}.
\end{align}
\end{subequations}
For the young households, the discount factor $B_{t+1}^i$ is itself a function of $\tau_t$ through $s_t$, so the factor $(1+B_{t+1}^i)$ in $\upsilon_{1,t}^i$ cannot be held fixed while differentiating. It is convenient to absorb it into the consumption level, defining
\begin{align}\label{\refeq:hatc1i}
    \hat{c}_{1,t}^i \equiv \left(1+B_{t+1}^i\right)^{\frac{1}{1-1/\rho}}\tilde{c}_{1,t}^i \qquad\Longrightarrow\qquad \upsilon_{1,t}^i = \frac{\left(\hat{c}_{1,t}^i\right)^{1-1/\rho}}{1-1/\rho},
\end{align}
so that a single numerical derivative of $\ln(\hat{c}_{1,t}^i)$ picks up both channels at once.\footnote{In implementation $\hat{c}_{1,t}^i$ is never formed as a level. Only two things are ever needed from it, and both avoid the $1/(1-1/\rho)$ power: $\left(\hat{c}_{1,t}^i\right)^{1-1/\rho} = \left(1+B_{t+1}^i\right)\left(\tilde{c}_{1,t}^i\right)^{1-1/\rho}$ and $\ln\left(\hat{c}_{1,t}^i\right) = \ln\left(1+B_{t+1}^i\right)/(1-1/\rho)+\ln\left(\tilde{c}_{1,t}^i\right)$. This matters as $\rho\rightarrow 1$: the literal definition raises $1+B_{t+1}^i$ to a power that diverges, overflowing in double precision well before $\rho$ reaches $1$ (at $\rho = 1.001$ the exponent is already $1001$), whereas neither expression above does.}

The two log-derivatives are treated differently. That of $\hat{c}_{1,t}^i$ --- and, as an input to it, that of $h_t$ --- is approximated numerically along the $\tau_t$ dimension of the solution grid, since it depends on $\tau_t$ through the continuation policies and has no closed form. By contrast, $d\ln(c_{2,t}^i)/d\tau_t$ \textbf{must} be taken from its closed form, given in \eqref{\refmodeleq:PEE}. Differentiating it numerically along the grid would be wrong, not merely imprecise: $s_{t-1,i}/s_{t-1}$ varies along that grid, and the policy maker takes it as predetermined, so the grid derivative would include a channel that does not belong in the first order condition. The closed form is exactly the expression already used in the log case, \eqref{\refmodeleq:PEELOG:v2i}, with one substitution: the analytical $\partial\ln(\Theta_{h,t})/\partial\tau_t$ is replaced by the numerically obtained $d\ln(h_t)/d\tau_t$. Since $h_t = \Theta_{h,t}(s_{t-1}/\nu_t)^{\alpha\xi/(1+\alpha\xi)}$ and $s_{t-1}$ is held fixed under $\partial/\partial\tau_t$, the two coincide, and the log and CRRA cases share one formula rather than two.

\smalltitle{The terminal period, and the log limit.}
The terminal period needs no numerical differentiation: by the chain rule, each term of \eqref{\refmodeleq:terminalPEECRRA} is a consumption level raised to $1-1/\rho$ times the \emph{same} log-derivative already available in closed form from the log case. What is genuinely new at $t=T$ is that $B_T^i$ is no longer the primitive $\beta_i$, so the terminal problem is $z_T(\tau_T, s_{T-1})$ --- state-dependent, unlike its log counterpart. The recursion at $t<T$, whose expressions divide by $1-1/\rho$, refuses $\rho = 1$ outright and directs the caller to \cref{\refsec:PEELOG} rather than approximating the limit. Since the log case here is solved by a wholly different route --- $T$ independent scalar problems with no state grid --- the two specifications provide a genuine cross-check on each other: solving the CRRA recursion at $\rho$ close to one must reproduce the log path, with any discrepancy that does not shrink as $\rho\rightarrow1$ indicating a defect in one of the two implementations rather than a property of the model. This check is cheap and worth keeping, as the two share almost no code.

\begin{alg}[Iterative grid search with boundary and feasibility checks]
\label{\refalg:CRRA:grid}
Fix $\bm{\theta}$ and global grids. The following identifies a sequence of policy functions $\tau_t(s_{t-1})$ as well as state policy functions $s_t(\tau_t, s_{t-1})$ that determine the economic state as a function of current taxes and the state along a politico-economic equilibrium path, iterating backwards from the terminal period $t=T$.

\smallskip\noindent\emph{Terminal period.}
Every object entering $z_T$ is closed form in $(\tau_T, s_{T-1})$. On the full grid $\mathcal{S}\tilde{\mathcal{T}}$, compute and store: $d \ln(h_t)/d \tau_t$ (closed form for $t=T$), $\Theta_{h,t}$ (closed form for $t=T$), $h_t$, $B_t^i$, $\Gamma_{s,t-1}$, $s_{t-1,i}/s_{t-1}$, $c_{2,t}^i$, and $\tilde{c}_{1,t}^i$ (closed form for $t=T$). Assemble the first order condition $z_T$ from these, and apply the selection rule of \cref{\refsec:RobustRoot} along $\tau$ for each state --- boundary checks, direction filter and objective integration included --- to obtain $\tau_T(s_{T-1})$ on $\mathcal{S}$. Finally, re-evaluate the auxiliary objects at the selected taxes and construct the interpolants $\tau^T(\cdot)$ and $h^T(\cdot)$ that the previous period will call.

\smallskip\noindent\emph{Iterating backwards for $t<T$.} Given the interpolants $\tau^{t+1}(\cdot)$ and $h^{t+1}(\cdot)$, each period proceeds in four steps.

\begin{enumerate}
    \item \emph{State approximation.} On the Cartesian grid $\tilde{\mathcal{T}}\times\mathcal{S}^{\prime}$ evaluate the forward pass \eqref{\refeq:stateApprox}, retaining $\Theta_{s,t}$ (and $\Theta_{h,t}$, $\Gamma_{s,t}$, $\tau^{t+1}$, $h^{t+1}$, $\bm{B}_{t+1}$ for later use). Form the residual \eqref{\refeq:stateResidual} for every $s_{t-1}\in\mathcal{S}$ by broadcasting, and locate its root along the $s_t$ dimension, giving $s_t(\tau_t, s_{t-1})$ on $\tilde{\mathcal{T}}\times\mathcal{S}$ together with the feasibility mask. This is the only step that touches $\mathcal{S}^{\prime}$, and the only one that calls the continuation interpolants.

    \item \emph{Current-period objects.} With $s_t$ resolved, evaluate on $\tilde{\mathcal{T}}\times\mathcal{S}$: $h_t\left(\Theta_{h,t}, s_{t-1}\right)$, then $R_t$, $\bm{B}_t$, $\Gamma_{s,t-1}\left(\bm{B}_t, \tau_t, \theta_t\right)$ and $s_{t-1,i}/s_{t-1}$ --- the last two evaluated at the previous period's parameter vintage, exactly as in the log case. Then the consumption levels entering \eqref{\refeq:PEEcrraTerms}: $\hat{c}_{1,t}^i$ from \eqref{\refeq:hatc1i} and $c_{2,t}^i$.

    \item \emph{Derivatives and the first order condition.} Approximate $d\ln(h_t)/d\tau_t$ and $d\ln(\hat{c}_{1,t}^i)/d\tau_t$ numerically along $\tau_t$, separately for each $s_{t-1}$; take $d\ln(c_{2,t}^i)/d\tau_t$ from its closed form as described above. Assemble $z_t$ on $\tilde{\mathcal{T}}\times\mathcal{S}$ via \eqref{\refeq:PEEcrraTerms} and the political weights of \eqref{\refeq:LOG:foc}, leaving infeasible cells undefined.

    \item \emph{Selection and reporting.} Apply the selection rule \eqref{\refeq:candidates} along $\tau_t$ for each $s_{t-1}$, skipping infeasible cells, to obtain $\tau_t(s_{t-1})$. Re-evaluate steps 1--2 at the selected taxes to obtain the solution dictionary on $\mathcal{S}$, and construct the interpolants $\tau_t(s_{t-1})$ and $h_t(s_{t-1})$ that the previous period will call. A light smoothing of the selected $\tau_t(s_{t-1})$ before interpolation removes small kinks inherited from the interpolated continuation policies; the smoother's interior knots are fixed in advance (at every $m$-th valid node) rather than selected from the data, so that the smoothed profile is a \emph{linear} map of the profile it smooths --- this keeps the composed map from parameters to solution continuous, which matters because the calibration of \cref{\refsec:calibration} differentiates through it.
\end{enumerate}

Steps 1--4 are stated for the whole grid $\mathcal{S}$ at once rather than for one $s_{t-1}$ at a time. Nothing forces this --- the states are independent and could equally be looped over --- but the grid evaluations dominate the cost only through the number of \emph{calls}, not the number of gridpoints, so handling all states in one pass is materially cheaper than looping, at the price of carrying the feasibility mask explicitly rather than subsetting the grid per state.
\end{alg}

\subsection{Calibration}\label{\refsec:calibration}
The calibration procedure is essentially a nested fixed point procedure. In the inner part, we solve for the full politico-economic equilibrium path. The outer part then solves the following system of equations (repeated from \cref{\refmodelsec:calibration}):
\begin{subequations}\label{\refeq:calibration}
    \begin{align}
        \hat{R}_{t_0} &= \alpha\left(\frac{s_{t_0-1}}{\nu_{t_0}h_{t_0}}\right)^{\alpha-1} \label{\refeq:calibration:R}\\
        \hat{\tau}_{t_0} &= \tau_{t_0} \label{\refeq:calibration:tau}
    \end{align}
\end{subequations}
where the ``hats'' reflect exogenous targets and $t_0$ indicates a specific baseline year. In our main specification, we assume that $\beta_{t,i} = \beta$ is constant over time and across types. Similarly, we assume that most exogenous parameters do not change over time except for the ageing population reflected by $\nu_t$. We adjust $\beta$ to target the interest rate and $\omega$ for the tax rate.

The full calibration procedure proceeds as follows: in an inner loop, given a set of parameter values, we first update all auxiliary parameter definitions to make sure that derived parameters are up to date, and then resolve the entire PEE path. An outer loop searches the parameter space for $\beta, \omega$ to ensure that \eqref{\refeq:calibration} holds.

\smalltitle{What the outer loop does not have to carry.}
Two features make this outer problem smaller than the corresponding one for the models with an informal sector, where four parameters have to be solved jointly with the path.

First, the pension characteristic $\theta$ is available in closed form from the replacement-rate ratio, \eqref{\refmodeleq:calibration:theta}, and depends only on the relative-income data $z_i^{\eta}$ --- not on $\Theta_{h,t_0}$ or any other equilibrium object. It is therefore computed once, before the loop starts, and never revisited. The same is true of $\bm{\eta}_i$ and $\bm{X}_i$: the eigenvector step of \cref{\refmodelsec:calibration:etaX} uses only survey data and $\xi$.

Second, in variant B the parameter $X$ is not a third dimension of the search. By \eqref{\refmodeleq:calibration:yxCommonX} it enters no aggregate, so neither target in \eqref{\refeq:calibration} contains it; the loop runs without it and $X$ is recovered afterwards from the closed form \eqref{\refmodeleq:calibration:Xsolve}, with no re-solve. It is nevertheless worth verifying numerically that re-solving at the recovered $X$ leaves $\beta$, $\omega$, $\tau_{t_0}$ and $R_{t_0}$ bit-for-bit unchanged: the claim is an identity, not an approximation, so any drift at all is a defect --- most likely an $\eta_i$ or $X_i$ that has leaked into an aggregate expression through something other than $y_i^{\eta}$.

\smalltitle{The interest rate target needs the path, not the baseline year.}
Note that \eqref{\refeq:calibration:R} involves $s_{t_0-1}$, the savings state \emph{entering} the baseline year. It is not an observable and not a free parameter: it is produced by the equilibrium path from the initial condition onwards, so the inner solve cannot be shortened to the baseline period. This is what makes $\beta$ a genuinely global object in the calibration --- it moves $s_{t_0-1}$ through every preceding period as well as $h_{t_0}$ within the baseline year --- and it is the reason the two targets in \eqref{\refeq:calibration} cannot be assigned to $\beta$ and $\omega$ one at a time and solved separately.

\smalltitle{Cost of the inner solve.}
With log preferences the inner solve is cheap: by \eqref{\refmodeleq:PEELOG:decoupling} the PEE path is $T$ independent scalar problems whose solution does not depend on $\beta$ through any continuation policy, only through the period-$t$ weights in \eqref{\refeq:LOG:foc}, after which the economic equilibrium follows by the forward pass of \cref{\refsec:EE}. Under CRRA preferences the inner solve is the full backward recursion of Algorithm \ref{\refalg:CRRA:grid}, and the state grid $\mathcal{S}$ has to be rebuilt whenever a trial $\beta$ moves the steady state materially; a grid fixed once at the initial guess will silently degrade as the outer loop wanders.

\subsection{Endogenous system characteristics}\label{\refsec:ESC}
This section documents how the design choices of \cref{\refmodelsec:ESC} are solved. Two solvers are described, one per timing --- the sequential choice needs none, since \eqref{\refmodeleq:esc:seqFOC} is a scalar first order condition of the same form as the tax one --- together with the calibration of the cost parameter and the checks that establish the structural claims rather than assuming them.

Throughout, $\Theta = \left\lbrace \theta^{(1)},\dots,\theta^{(M_{\theta})}\right\rbrace$ denotes a grid on $[0,1]$ and $\mathcal{T}_t(\theta)$ the ordinary equilibrium tax at date $t$ given a design $\theta$, i.e.\ the solution of $z_t(\tau,\theta)=0$ selected as in \eqref{\refeq:candidates}. By the decoupling property this is a \emph{static} scalar problem, one root per node, with no dependence on $s_{t-1}$ and no recursion --- which is what makes the log case cheap.

\smalltitle{Why a grid and not a first order condition.} The question being asked is whether the choice is interior. A solver built on $\mathrm{d}\mathcal{W}/\mathrm{d}\theta = 0$ answers it only by failing to converge, and a corner is then indistinguishable from a badly behaved residual. We therefore evaluate the objective itself on $\Theta$, take the maximiser, and refine it by fitting a parabola through the three nodes around it when the maximum is interior. A corner is returned \emph{as} a corner. The cost is one objective evaluation per node; the benefit is that the central result of the exercise is observed rather than inferred. This inverts the convention used for the tax: there the first order condition is the cheap object and the objective is reconstructed from it (\cref{\refsec:RobustRoot}); here the objective is evaluated directly, precisely because its corner behaviour is the substance of the question.

\smalltitle{Evaluating the objective does not relax the discipline on predetermined states.} Direct evaluation may look like it surrenders what the first order condition was built to guarantee --- that the electorate takes the inherited distribution of savings as given. It does not, because that guarantee never resided in the representation. The tax residual $z_t$ excludes $s_{t-1,i}/s_{t-1}$ from the variation because it is \emph{constructed} as the derivative holding the ratio fixed; the grid implements the same definition by evaluating every candidate at the same predetermined ratio, and the two representations enforce one discipline. Under the leaded timing the discipline is automatic: the candidate $\theta_{t+1}$ does not appear in \eqref{\refmodeleq:esc:auxiliary:si} at vintage $t-1$, so no channel exists through which the sweep could move the ratio. What \emph{does} move with the candidate --- the current young's labour, savings and their distribution --- is meant to: those choices are being made at $t$, not sunk. Under the permanent timing the candidate does enter the ratio's formula, through $\theta_{t_0}$, and there the ratio is pinned explicitly; \cref{\refsec:ESC:permanent} states the pinning and measures what ignoring it would cost.


\subsubsection{The leaded choice}\label{\refsec:ESC:leaded}
The design chosen at $t$ is a state at $t+1$, so the object to identify is a pair of sequences of policy functions, $\mathcal{T}_t(\theta_t)$ and $\Theta_t(\theta_t)$. \Cref{prop:esc:separability} says the second is constant in its argument, but the algorithm does not use that: it tabulates the policy on $\Theta$ precisely so that the claim can be measured.

\begin{alg}\label{\refalg:esc:leaded}
\emph{Leaded choice, log preferences.} Iterate backwards over $t = T-1,\dots,1$.
\begin{enumerate}
    \item \emph{Terminal period.} Solve $\mathcal{T}_{T-1}(\theta)$ for every $\theta\in\Theta$ from the terminal first order condition, i.e.\ \eqref{\refmodeleq:PEELOG} evaluated with $\beta_{T,j}=0$. There is no $\theta_T$ to choose.
    \item \emph{Taxes.} At $t<T-1$, evaluate $z_t$ on the flattened mesh $\Theta\times\tilde{\mathcal{T}}$ in a single call --- $z_t$ is elementwise in both arguments, so the mesh costs one evaluation rather than $M_{\theta}$ --- and select the maximiser at each node as in \eqref{\refeq:candidates}. \emph{Polish} each interior selection by a bracketed scalar root find within one grid cell.
    \item \emph{Design.} For each $\theta_t\in\Theta$, evaluate
    \begin{align}\label{\refeq:esc:leadedObjective}
        \mathcal{W}_t\left(\theta';\theta_t\right) = \sum_i \gamma_{t-1,i}\omega p_{t-1}\mu_{t-1,i}\ln\left(c_{2,t}^i\right) + \nu_t\sum_i \gamma_{t,i}\mu_{t,i}\left[(1+\beta_{t,i})\ln\left(\tilde{c}_{1,t}^i\right)+\beta_{t,i}\ln\left(R_{t+1}\right)\right]
    \end{align}
    at every candidate $\theta'\in\Theta$, using $\tau_t = \mathcal{T}_t(\theta_t)$ from step 2 and the continuation $\tau_{t+1} = \mathcal{T}_{t+1}(\theta')$, $\theta_{t+2} = \Theta_{t+1}(\theta')$, $\tau_{t+2} = \mathcal{T}_{t+2}(\theta_{t+2})$ from the periods already solved. Set $\Theta_t(\theta_t)$ to the refined maximiser.
\end{enumerate}
\end{alg}

Three implementation points are worth recording, because each guards against a plausible wrong answer.

\emph{The polish is not optional.} The grid selection alone is $O(\text{spacing}^2)$ accurate, roughly $10^{-4}$ at the resolutions used, while $\tau_{t_0}$ is a calibration target matched at $10^{-8}$. Without the refinement the cost parameter would be calibrated against grid noise. Corners are never polished: $z_t \neq 0$ there is the correct answer, not a residual.

\emph{$s_{t-1}$ is set to one.} By \cref{prop:esc:separability} the maximiser does not depend on it. The normalisation is convenient rather than necessary and is verified numerically, by re-maximising \eqref{\refeq:esc:leadedObjective} at savings levels spanning two orders of magnitude and confirming the argument does not move.

\emph{Interpolation is piecewise linear throughout.} A policy function that is about to be maximised over must not carry interpolation artefacts: a smoother or a cubic can manufacture a spurious interior maximum out of a profile that is genuinely flat at a bound, and an argmax --- unlike a root --- has no residual through which such an artefact would announce itself. The cautions of \cref{\refsec:PEE} about interpolation kind therefore apply with more force here than in their original setting.


\subsubsection{The permanent choice}\label{\refsec:ESC:permanent}
Here the design is chosen once, at $t_0$, and the electorate chooses $\tau_{t_0}$ alongside it. \Cref{prop:esc:concentration} reduces what looks like a two-dimensional grid search over $\left(\tau_{t_0},\theta\right)$ to a one-dimensional one, and the reduction is what makes the exercise cheap:

\begin{alg}\label{\refalg:esc:permanent}
\emph{Permanent choice.} Given a pinned ratio $\bar{\sigma}_i$, for each $\theta\in\Theta$: (i) obtain $\tau_t = \mathcal{T}_t(\theta)$ for every $t\geq t_0$ --- for $t>t_0$ this is the ordinary equilibrium at a constant design, with no recursion in $\theta$ to solve; (ii) evaluate \eqref{\refeq:esc:leadedObjective} at $t_0$ with $\theta_{t_0}=\theta'=\theta$ and $s_{t_0-1,i}/s_{t_0-1}=\bar{\sigma}_i$. Take the refined maximiser over $\Theta$. The vote being anticipated, $\bar{\sigma}_i$ is itself the ratio implied by the winning design, so the design is the fixed point of
\begin{align}\label{\refeq:esc:permFixedPoint}
    \theta^{\ast} = \arg\max_{\theta\in\Theta}\ \mathcal{W}_{t_0}\left(\theta;\ \sigma_i(\theta^{\ast})\right),
\end{align}
reached by iterating the above from the incumbent design. Convergence is reported, not assumed: this is a best response and not a contraction, so a cycle has to be visible.
\end{alg}

\smalltitle{The predetermined ratio must be held fixed, and at which value.} This is the one place where the permanent timing differs from the leaded one in substance rather than in bookkeeping, and it raises two questions that are easily run together.

The first is whether the ratio may move with the candidate at all. It may not. Under the leaded timing $\theta_{t+1}$ does not appear in $s_{t_0-1,i}/s_{t_0-1}$ and the question never arises; under the permanent timing $\theta$ does appear there, through $\theta_{t_0}$. But savings made at $t_0-1$ are sunk when the vote is taken at $t_0$: a voter comparing two candidates does not face a different $s_{t_0-1,i}$ under each. Maximising $\mathcal{W}_{t_0}$ while recomputing the ratio at every candidate would credit the electorate with internalising a state it takes as given --- the same error that \cref{\refsec:PEE} forbids for $\tau_t$, arriving in a new guise --- and it is not a small one: at $p=0.4$ the two readings give $0.775$ and $0.910$. It is also inconsistent with the tax it would be paired with, since $s_{t_0-1,i}/s_{t_0-1}$ is a function of $\tau_{t_0}$ as well as $\theta_{t_0}$ while $\mathcal{T}_{t_0}$ solves $z_{t_0}=0$, which is constructed holding the ratio fixed. It would take \cref{prop:esc:concentration} with it.

The second question is which value the pinned ratio takes, and there the answer is the timing's. The reform is \emph{anticipated}: households arrive at $t_0$ knowing a design will be chosen, and what they choose among is which of the $\left|\Theta\right|$ equilibrium paths the economy follows from $t_0$ on. The savings they made at $t_0-1$ were therefore made against the design that wins, which is \eqref{\refeq:esc:permFixedPoint}. The incumbent's ratio --- the natural pinning if the reform were a surprise --- is the seed of that iteration and not its answer. The two coincide \emph{exactly} wherever the chosen design reproduces the incumbent one, which is precisely the calibration condition \eqref{\refeq:esc:calibration}, so the calibrated cost is common to both readings and only the counterfactuals separate them; at $p=0.4$, away from the permanent timing's own calibration, the anticipated reading gives $0.775$ against the unanticipated $0.773$, and at $p=0.25$ it gives $0.542$ against $0.549$. Both are reported.

That the fixed point is a scalar problem at all is worth one line. The ratio in \eqref{\refmodeleq:esc:auxiliary:si} at vintage $t_0-1$ is a function of date-$t_0$ policy alone, so the kinked path the reform actually is --- the incumbent design before $t_0$, the chosen one from $t_0$ --- delivers the same $\bar{\sigma}_i$ as a design that had always been $\theta$, and no separate solve is needed to evaluate a candidate's savings. The kink is nonetheless carried in the reported equilibrium path, where a design applied to periods before the vote would be a counterfactual past.

\smalltitle{Verification.} \Cref{prop:esc:concentration} is checked rather than trusted: the tax at the chosen design must equal $\mathcal{T}_{t_0}$ evaluated there, and does, to $10^{-6}$. Were the concentration to fail --- because some channel made $\theta$ a function of $\tau$ --- \cref{\refalg:esc:permanent} would be solving the wrong problem, and this is the check that would say so.


\subsubsection{CRRA preferences}\label{\refsec:ESC:crra}
None of the log simplifications survive. The powers in $\left(\tilde{c}\right)^{1-1/\rho}$ do not turn products into sums, so \eqref{\refmodeleq:esc:separability} fails, both $s_{t-1}$ and $\theta_t$ re-enter the design choice, and the Markov object is a policy function over the two-dimensional state $\left(s_{t-1},\theta_t\right)$ --- a second dimension on top of the grid machinery of \cref{\refsec:PEE}.

For the permanent timing this costs nothing: a permanent design means a \emph{constant} $\theta$ path, each candidate is one ordinary CRRA equilibrium solve, and \cref{\refalg:esc:permanent} goes through unchanged. For the leaded timing we solve the equilibrium \emph{path} rather than the policy function:

\begin{alg}\label{\refalg:esc:crraPath}
\emph{Leaded choice, CRRA preferences.} Given a design path $\bm{\theta}$ with $\theta_t$ fixed at the incumbent value for $t\leq t_0$: sweep forward over $t \geq t_0$; at each $t$, evaluate the CRRA analogue of \eqref{\refeq:esc:leadedObjective} at every candidate $\theta_{t+1}\in\Theta$, re-solving the whole politico-economic equilibrium for each candidate while holding $\theta_{t+2},\dots$ at the current iterate; replace $\theta_{t+1}$ by the maximiser. Repeat the sweep until the path stops moving.
\end{alg}

Re-solving at every candidate is what makes the envelope logic right: $\tau_t$ is re-optimised at each design, so its own response to $\theta_{t+1}$ contributes nothing to first order and the objective is evaluated at the equilibrium tax rather than a stale one.

\smalltitle{What \cref{\refalg:esc:crraPath} assumes, and how far it is true.} Holding $\theta_{t+2}$ fixed while $\theta_{t+1}$ varies is exact if the choice at $t+1$ does not respond to the design it inherits. Under log preferences that is \cref{prop:esc:separability} and no approximation at all. Under CRRA it is an approximation, and we measure it directly: perturbing $\theta_{t+1}$ and re-optimising $\theta_{t+2}$ at each perturbation gives $\mathrm{d}\theta_{t+2}/\mathrm{d}\theta_{t+1}$ of $-0.0085$ at $\rho=2$. The residual state dependence is two orders of magnitude below the movements being studied, so the path iteration is adequate here; at a calibration where it is not, the two-dimensional grid becomes necessary and this shortcut has to go.

\smalltitle{Validation against the log limit.} The CRRA solver has no closed form to check against, so it is held to continuity: as $\rho\rightarrow1$ its answer must converge on the log solver's, and at a first-order rate, since nothing in the problem is kinked at $\rho=1$. Holding the cost parameter and the calibration fixed, the gap is $0.0903$, $0.0448$ and $0.0188$ at $\rho = 1.10$, $1.05$ and $1.02$, so the ratio $\text{gap}/(\rho-1)$ is $0.90$, $0.90$, $0.94$. The two solvers share the weights of \eqref{\refeq:esc:leadedObjective} and nothing else --- one maximises a closed form on a state grid by backward recursion, the other iterates on a path of full equilibrium solves --- so their agreement in the limit is informative.

\smalltitle{The exact solution: policy functions over the two-dimensional state.} \Cref{\refalg:esc:crraPath} approximates the Markov object; the object itself is also computed, by extending the backward recursion of \cref{\refsec:PEE} with the design as a second state:

\begin{alg}\label{\refalg:esc:crra2D}
\emph{Leaded choice, CRRA preferences, exact.} Let $\mathcal{S}\times\Theta$ be a grid over the state $\left(s_{t-1},\theta_t\right)$ and $\Theta'$ a grid of candidate designs. Solve the terminal period as in \cref{\refsec:PEE} once per $\theta_T\in\Theta$. Then, backwards over $t$:
\begin{enumerate}
    \item For every $\left(\tau_t, s_{t-1}, \theta_{t+1}\right) \in \mathcal{T}\times\mathcal{S}\times\Theta'$, resolve the equilibrium state $s_t$ from the state residual of \cref{\refsec:PEE}, with the continuation $\left(\tau_{t+1}, h_{t+1}\right)$ read off the period-$(t+1)$ policy interpolants at $\left(s_t, \theta_{t+1}\right)$.
    \item Assemble $z_t$ on that grid. The inherited design $\theta_t$ enters only the current old's $\mathrm{d}\upsilon_{2,t}^i/\mathrm{d}\tau_t$ term --- it splits benefits already granted --- so the numerical $\tau$-derivatives are shared across the $\theta_t$ axis and only that one term is recomputed per node.
    \item Select $\tau^{\ast}\left(s_{t-1},\theta_t,\theta_{t+1}\right)$ by the downward-crossing criterion of \cref{\refsec:RobustRoot}, evaluate the objective at the selection, and take $\Theta_t\left(s_{t-1},\theta_t\right)$ as its maximiser over $\Theta'$, refined parabolically when interior; corners are reported as corners.
    \item Tabulate $\left(\tau_t, \theta_{t+1}, s_t, h_t\right)$ at the chosen policies and rebuild the interpolants for period $t-1$.
\end{enumerate}
\end{alg}

One backward pass --- no iteration and no convergence tolerance; periods at which the design is history rather than a choice (every $t\leq t_0$ in the baseline timing) are solved with the candidate set collapsed to the inherited design, which matters here because $\tau_t$ genuinely responds to $\theta_{t+1}$ under CRRA. Pinned everywhere, the recursion collapses to the exogenous-design solver of \cref{\refsec:PEE} on a redundant grid, and reproduces its tax path to $\max_t\left|\Delta\tau_t\right| = 3.2\times10^{-5}$ --- a check of every ingredient except the choice layer against the production solver. At the calibrated cost at $\rho=2$, the exact choice at $t_0$ is $0.759$ against the path iteration's $0.761$, a gap of $0.002$, and the tax at $t_0$ still lands on its target (drifts of $3.2\times10^{-4}$ in $\tau$ and $3.1\times10^{-4}$ in $R$). The objective is flat enough near its maximum that either method identifies the design only to about $\pm 0.01$, and the two agree comfortably within that band. The discretisation requirements are properties of the grids rather than of the cost parameter: the state grid is immaterial (the choice moves by $4\times10^{-4}$ between $50$ and $150$ nodes) while the design-\emph{state} grid is not ($7$ nodes misplace the choice by $0.02$; $13$ suffice). \Cref{\refalg:esc:crraPath} is therefore confirmed adequate for the counterfactual tables, with \cref{\refalg:esc:crra2D} as the standard it is held to.


\subsubsection{Calibrating the cost}\label{\refsec:ESC:calibration}
Three targets and three parameters. $\left(\beta,\omega\right)$ are calibrated to $\left(R_{t_0},\tau_{t_0}\right)$ exactly as in \cref{\refsec:calibration}, with the cost installed and the design held at $\theta^{\ast}$; $\phi$ is imposed; and $p$ is chosen so that the design \emph{in force} at the baseline year, on a path where the political choice binds from the first period of the horizon, is the observed one,
\begin{align}\label{\refeq:esc:calibration}
    \theta_{t_0} = \Theta_{t_0-1}\left(\theta_{t_0-1}\right) = \theta^{\ast}.
\end{align}
Read in words: the model's own political history delivers the observed system by the baseline year. Everything after $t_0$ is then a prediction, and the drift of the design as $\nu_t$ falls is the object of interest.

\smalltitle{Why the design in force and not the choice made.} The design is a state chosen one period earlier, so $\theta_{t_0}$ and $\Theta_{t_0}(\theta^{\ast})$ are different numbers: the first was picked under $\nu_{t_0-1}$, the second binds only at $t_0+1$. They would coincide if the policy function were time invariant, and $\nu_t$ is not. The distinction is forced by what the counterfactuals report. Those are equilibrium paths read at $t_0$ (\cref{\refsec:ESC:counterfactuals}), so the baseline path must sit on the observed design \emph{at} $t_0$ or every comparison in the tables is drawn against a number the model missed. Under log preferences \cref{prop:esc:separability} makes $\Theta_{t_0-1}$ constant in its argument, so \eqref{\refeq:esc:calibration} does not depend on the inherited $\theta_{t_0-1}$ and no fixed point in the design's own history arises; under CRRA the residual is evaluated with the rest of the path held at $\theta^{\ast}$, exactly as elsewhere in this section, which is again exact at the root. The two targets are close --- calibrating on the choice at $t_0$ instead gives $p = 0.402$ against $0.408$ under the proportional specification at $\rho=1$, and a design at $t_0$ of $0.727$ against the observed $0.738$ --- so what the choice of target buys is not a different economics but a baseline row that reproduces its own datum.

\smalltitle{The nesting is an approximation, and is reported as one.} $\left(\beta,\omega\right)$ are calibrated along the path with the design held fixed, not along the endogenous-design path. At the calibrated $p$ the two agree at $t_0$ and $t_0+1$ by construction and drift only later, so the error is second order --- but it is an error, and the solved endogenous path is therefore checked back against the $\left(R,\tau\right)$ targets it was calibrated to. The residuals are $9\times10^{-15}$ and $2.6\times10^{-9}$, which is what licenses treating the two blocks as separable.

This caveat belongs to the leaded timing alone. Under the permanent timing the chosen design is constant from $t_0$ on and, at the calibrated $p$, equal to the incumbent one --- so the exogenous path $\left(\beta,\omega\right)$ are calibrated along \emph{is} the endogenous path, period for period, and the nesting is an identity rather than an approximation.

\smalltitle{The permanent residual is the fixed-point condition.} Equation \eqref{\refeq:esc:calibration} asks for the $p$ at which the electorate re-elects the system it has, and under the permanent timing that is exactly \eqref{\refeq:esc:permFixedPoint}: a cost at which the choice taken against the incumbent's savings ratio reproduces the incumbent design is a cost at which the incumbent design is the fixed point. The two therefore share a root, and the residual is evaluated in the cheaper single-pass form; only away from the root do they differ, so a scan must not mix them.

\smalltitle{Why \eqref{\refeq:esc:calibration} is scanned and not bracketed.} The residual is \emph{flat at a corner}: wherever the choice sits at $\theta=0$ or $\theta=1$ it does not move with $p$ at all. A root finder handed such a bracket reports spurious convergence at whichever endpoint it happened to evaluate. We therefore scan $p$ on a logarithmic grid, locate a genuine sign change, and only then bracket. A run that finds no interior crossing reports that fact rather than returning a number --- which is the correct outcome for France, whose observed design is itself at the boundary. It is also what happens if the scanned interval is simply placed wrong: the required cost spans an order of magnitude across $\rho$, so an interval chosen at $\rho=2$ lies entirely below the root at $\rho=0.5$, and the scan says so instead of converging on its own endpoint. The interval scanned therefore spans all three.

\smalltitle{Two consequences of the identification.} Under the flat-only specification, \eqref{\refmodeleq:esc:RR} makes $\theta^{\ast}$ a function of $p$, so the target in \eqref{\refeq:esc:calibration} moves with the parameter being calibrated. This is handled by re-deriving $\theta^{\ast}$ from $RR_0$ at every trial $p$ --- the datum is the replacement-rate ratio, not the design --- and the inversion is bracketed by a scan, since $(1-\theta)/\theta$ is decreasing while $f$ is increasing and their product need not be monotone. Second, the required cost turns out to be nearly the same under all three timings, which is itself a useful check that the two solvers are solving the same economic problem.

\smalltitle{The calibrated cost across $\rho$.} Under CRRA the residual in \eqref{\refeq:esc:calibration} is evaluated with the future design held at $\theta^{\ast}$ while the single choice at $t_0$ varies --- exact at the root, where the held path \emph{is} the equilibrium path, and only an approximation away from it, where it merely has to locate the sign change. The calibrated cost falls steeply in the intertemporal elasticity: under the proportional specification $p = 0.965,\ 0.408,\ 0.090$ at $\rho = 0.5,\ 1,\ 2$ (and $1.461,\ 0.709,\ 0.178$ under the flat-only one, where $\theta^{\ast}$ moves with it: $0.690,\ 0.716,\ 0.733$). A higher elasticity strengthens the young's forward-looking stake in a Bismarckian design --- the same ordering the stake decomposition found --- so less friction is needed to keep the choice off the Beveridgean corner. All six points are in the repository's \texttt{results/esc/escCalibration\{,CRRA\}.csv}, produced by \texttt{python/US/runESC.py} and \texttt{runESCcrra.py} and consumed by the paper pipeline's \texttt{US\_ESC\_Calibration} table.


\subsubsection{The counterfactuals}\label{\refsec:ESC:counterfactuals}
Every counterfactual in this appendix --- and every one in the main text's US tables --- is a \emph{separate equilibrium path}, not a shock applied to the calibrated economy at the baseline year. The changed characteristics hold over the whole horizon, the economy begins at its own steady state rather than at the United States', the political choice binds from the first period, and the path is read at $t_0$. What a row describes is a country that has always had this mix of characteristics --- mostly American, partly French --- and the question the tables put is how far the observable characteristics carry the United States towards France.

\smalltitle{Why not an unanticipated reform at $t_0$.} An alternative construction dates the change at $t_0$, seeds the economy with the savings the baseline generation actually made, and so describes an economy surprised in 2020. That object is coherent, but it is not commensurable with France's own calibrated path, which is an equilibrium and not a surprise. Since the point of the French experiments is a comparison across countries, the paths have to be built the same way on both sides. The new-path construction has a second advantage: with the design pinned as history through $t_0$, a change dated $t_0$ could not move $\theta$ before $t_0+1$, so a table read at $t_0$ would be identical across the exogenous- and endogenous-design readings by construction. Here $\theta_{t_0}$ is itself an outcome, and 2020 carries the design response.

\smalltitle{What the convention costs, measured.} At $\rho=1$ the choice between the two constructions moves the average workweek and nothing else: the tax rate and the savings rate come out identical to every printed digit under both readings. Both are rate objects, and under log preferences with a Cobb--Douglas technology neither depends on the inherited capital stock, while the \emph{level} of hours responds to the wage and hence to $k_{t_0}$. This is a property of the log case and does not survive it --- under CRRA the tax responds to the state, and the $\rho\neq1$ rows move.

\smalltitle{Two readings.} Each counterfactual is reported both with the design held at the value the replacement-rate data imply under the changed characteristics, and with the design chosen. The two differ already at $t_0$, which is what makes a table read there informative. Alongside them the French tables carry France's own calibrated path. That row is not a counterfactual on the American model: France brings its own political weight $\omega$ as well as its own characteristics, so the distance between it and the all-characteristics row measures what the observables leave unexplained. Its workweek is a calibration target rather than a prediction, since the French calibration of \cref{\refsec:calibration} targets French hours relative to American ones, and is marked as such wherever it is printed.

